\documentclass[12pt]{article}
\usepackage[utf8]{inputenc}
\usepackage[russian]{babel}
\usepackage{amsmath}
\usepackage{amssymb}
\usepackage{amsthm}
\usepackage{geometry}

\geometry{a4paper, left=20mm, right=20mm, top=20mm, bottom=20mm}

\newtheorem{definition}{Определение}
\newtheorem{remark}{Замечание}

\begin{document}

\section*{Задача 3. Касательное пространство к множеству $M$ в точке $(0,0)$}

\[
M = \{(x_1, x_2) \in \mathbb{R}^2 \mid x_1 \le 0,\ x_2 \le 0,\ x_1^3 + x_2^3 \le 1\}, \qquad p = (0, 0).
\]

\subsection*{Часть I. Теоретические основания}

\subsubsection*{1.1. Что такое касательный вектор?}

Согласно определению из лекции и Википедии:

\begin{definition}[Касательный вектор к подмногообразию]
\textbf{Касательный вектор} в точке $p$ гладкого подмногообразия $M$ евклидова пространства --- это \textbf{вектор скорости} $\alpha'(0)$ некоторой гладкой кривой $\alpha$, целиком лежащей в $M$ и проходящей через $p$:
\[
\alpha: (-\varepsilon, \varepsilon) \to M, \quad \alpha(0) = p.
\]
Две кривые $\alpha$ и $\beta$ задают один и тот же касательный вектор, если $\alpha'(0) = \beta'(0)$ (они \emph{касаются} друг друга в точке $p$).
\end{definition}

\begin{remark}
Для множеств с границей (как наше $M$) кривая может быть определена лишь на полуинтервале $[0, \varepsilon)$, так как из граничной точки нельзя двигаться ``наружу''.\\ 
\\
Поэтому мы рассматриваем кривые $\alpha: [0, \varepsilon) \to M$.
\end{remark}

\subsubsection*{1.2. Касательный конус и касательное пространство}

\begin{definition}[Касательный конус]
Множество всех возможных векторов скорости $\alpha'(0)$ по всем кривым $\alpha: [0, \varepsilon) \to M$ с $\alpha(0) = p$ называется \textbf{касательным конусом} (конусом допустимых направлений):
\[
K_M(p) = \left\{ v \in \mathbb{R}^n \;\middle|\; \exists\, \alpha: [0,\varepsilon) \to M,\ \alpha(0) = p,\ \alpha'(0) = v \right\}.
\]
\end{definition}

\begin{definition}[Касательное пространство]
\textbf{Касательное пространство} к $M$ в точке $p$ --- это линейная оболочка касательного конуса:
\[
T_p M = \operatorname{span} K_M(p) = \left\{ \sum_{i=1}^k \lambda_i v_i \;\middle|\; k \in \mathbb{N},\ \lambda_i \in \mathbb{R},\ v_i \in K_M(p) \right\}.
\]
Это наименьшее линейное подпространство $\mathbb{R}^n$, содержащее все касательные направления.
\end{definition}

\begin{remark}[]
Для открытого множества $U \subset \mathbb{R}^n$ касательное пространство $T_p U \cong \mathbb{R}^n$, так как из любой точки можно двигаться в любом направлении.\\
\\
Для множества с границей касательный конус $K_M(p)$ может быть строго меньше $\mathbb{R}^n$, но его линейная оболочка всё равно может дать всё $\mathbb{R}^n$ (как в нашей задаче).
\end{remark}

\subsection*{Часть II. Решение}

\textbf{1. Анализ ограничений.\\}

Введём $g_1 = x_1$, $g_2 = x_2$, $g_3 = x_1^3 + x_2^3 - 1$. В точке $p$:
\[
g_1(0,0) = 0, \quad g_2(0,0) = 0, \quad g_3(0,0) = -1 < 0.
\]

$\forall (x_1, x_2) \in M$: $x_1 \le 0, x_2 \le 0 \Rightarrow x_1^3 + x_2^3 \le 0 < 1$, поэтому $g_3$ неактивно на $M$.\\

\textbf{2. Необходимое условие.\\}

Пусть $\alpha: [0, \varepsilon) \to M$, $\alpha(0) = (0,0)$, $\alpha'(0) = v = (v_1, v_2)$.\\

$\alpha(t) \in M \Rightarrow \alpha_1(t) \le 0,\ \alpha_2(t) \le 0\ \forall t \in [0, \varepsilon)$.

\[
v_1 = \alpha_1'(0) = \lim_{t \to 0^+} \frac{\alpha_1(t)}{t} \le 0, \quad v_2 = \alpha_2'(0) = \lim_{t \to 0^+} \frac{\alpha_2(t)}{t} \le 0.
\]

\textbf{3. Достаточное условие.\\}

$\forall v = (v_1, v_2)$ с $v_1 \le 0, v_2 \le 0$ рассмотрим $\alpha(t) = (tv_1, tv_2)$, $t \in [0, \varepsilon)$:
\[
\alpha_1(t) = tv_1 \le 0, \quad \alpha_2(t) = tv_2 \le 0, \quad \alpha_1^3 + \alpha_2^3 = t^3(v_1^3 + v_2^3) \le 0 < 1.
\]

$\alpha(0) = p$, $\alpha'(0) = v \Rightarrow v \in K_M(p)$.\\

\textbf{4. Касательный конус.}

\[
K_M(0,0) = \{(v_1, v_2) \in \mathbb{R}^2 \mid v_1 \le 0,\ v_2 \le 0\} \quad \text{(третий квадрант)}.
\]

\textbf{5. Касательное пространство.\\}

$u_1 = (-1, 0) \in K_M(0,0)$, $u_2 = (0, -1) \in K_M(0,0)$.\\

$\{u_1, u_2\}$ --- базис $\mathbb{R}^2$, поэтому:
\[
T_{(0,0)}M = \operatorname{span} K_M(0,0) \supseteq \operatorname{span}\{u_1, u_2\} = \mathbb{R}^2.
\]

\subsection*{Ответ}

\[
\boxed{T_{(0,0)}M = \mathbb{R}^2.}
\]

\begin{remark}[]
Точка $(0,0)$ --- это ``угол'' множества $M$ (пересечение двух граничных прямых $x_1 = 0$ и $x_2 = 0$).\\

Касательный конус в этой точке --- третий квадрант (все направления, ``смотрящие внутрь'' множества).\\

Хотя конус не совпадает со всем $\mathbb{R}^2$, он содержит два линейно независимых направления $(-1,0)$ и $(0,-1)$, поэтому его линейная оболочка заполняет всю плоскость.
\end{remark}

\end{document}
